\section{质点运动学}
\sourceproblem{选择1}{1}{7}
\problemtext 判断一维运动何时趋向原点。
\workings 到原点距离为 $|x|$；在 $x\ne0$ 时，距离减小等价于 $x\dot x<0$，即 $\mathrm d(x^2)/\mathrm dt<0$，选 C。
\crosscheck $x>0$ 时须 $\dot x<0$，$x<0$ 时须 $\dot x>0$；在原点不能称为正在趋近原点。
\sourceproblem{选择2}{1}{7}
\problemtext 识别瞬时速率的等价表达式。
\workings $\bm v=\mathrm d\bm r/\mathrm dt$，$v=|\bm v|=\mathrm ds/\mathrm dt=\sqrt{\dot x^2+\dot y^2+\dot z^2}$，也等于速度坐标分量所构成矢量的模。故原式(2)(3)(4)(5)正确，选 B。
\crosscheck 速度矢量本身不是其大小；弧长 $s$ 为累计路程。
\sourceproblem{选择3}{1}{7}
\problemtext $x=5+3t^2-t^3$（SI），判断运动方向与速率变化。
\workings $v=3t(2-t)$，$a=6(1-t)$。$0<t<1$ 时二者均正；$1<t<2$ 时 $v>0,a<0$；$t>2$ 时二者均负。故选 B。
\crosscheck 速率增减由 $va$ 的符号确定，$t=2$ 是速度换向点。
\sourceproblem{选择4}{1}{7}
\problemtext $\bm r=(3t^2+2)\bm i+6t^2\bm j$，判别轨迹及速率。
\workings 消去 $t$ 得 $y=2(x-2)$，$x\ge2$；$\bm v=6t\bm i+12t\bm j$，速率为 $6\sqrt5|t|$。选 B（变速直线）。
\crosscheck 加速度 $6\bm i+12\bm j$ 恒定；轨迹为直线上的半直线。
\sourceproblem{选择5}{1}{7}
\problemtext $x=t^2-40t$，求 $t_1$ 至 $t_2$ 的平均速度。
\workings 当 $t_2>t_1$ 时，
\[
\bar v=\frac{(t_2^2-t_1^2)-40(t_2-t_1)}{t_2-t_1}=t_1+t_2-40.
\]
选 A。
\crosscheck $t_2\to t_1$ 时趋于瞬时速度 $2t_1-40$，不可把平均速度当平均速率。
